Simultaneous interpreting (SI) requires an interpreter to listen, analyze, retain information, and produce another language under severe temporal constraints. Segment-onset delay is consequently a natural timing signal. It is distinct from ear--voice span (EVS), which requires an explicit choice of source--target linguistic anchors. Yet a temporal offset alone does not answer a different question faced by users and professional evaluators: does an interpretation \emph{feel} delayed? A short onset gap can coexist with omissions, broken delivery, or an interpreter who has stopped keeping up; conversely, a longer but fluent and complete segment may be judged less harshly.

This distinction creates an underexplored NLP evaluation problem. We ask whether a system can predict a professionally rated \emph{promptness score}: a reverse-oriented judgment for which a larger value means \emph{less} experienced delay. The answer matters for feedback systems and future SI assistants because it connects observable timing and text signals to the aspect of promptness identified in professional evaluation.

We use three scope distinctions throughout. At the \emph{task} level, the outcome is a professional promptness rating rather than elapsed delay or validated real-time perception. At the \emph{measurement} level, the protocol specifies an equal-weight aggregate target from two comparably qualified evaluators; cross-rater commonality remains evaluator-dependent. At the \emph{generalization} level, source-speech-group-held-out evaluation tests unseen source content, whereas interpreter-disjoint evaluation tests output from a held-out interpreter while allowing source-speech overlap.

We compare three deployable signal families: segment-onset delay, deterministic lexical/structural features from the source and interpreted output, and automatically predicted professional quality. A frozen shared encoder with two rubric-supervised quality outputs predicts LQ and EXP; a fold-local Ridge model combines these estimates with timing and optional structural features. In source-speech-group-held-out evaluation, the held-out group is excluded from upstream quality training and checkpoint selection; inner partitions provide cross-fitted quality predictions for promptness-model training. Interpreter-disjoint evaluation excludes every segment from the held-out interpreter throughout quality estimation and promptness-model training.

The results support a signal-decomposition account. The primary comparison asks whether professionally supervised quality signals contain information absent from onset timing; it does not establish the neural system as the strongest deployable text-based model. Deterministic lexical/structural+delay provides that stricter baseline, with predicted quality adding a smaller, seed-sensitive association. Under interpreter-disjoint evaluation, the full system reaches $r=\LOIOFullPearson\pm\LOIOFullPearsonSD$, with substantial variation across interpreters. Cross-rater analyses identify positive but asymmetric associations, and professional comments illustrate how evaluators describe low promptness ratings despite short onset delay.

Our contributions are:
\begin{itemize}
    \item We formulate a reverse-oriented professional promptness-rating task distinct from segment-onset delay and define an aggregate target with evaluator-specific analysis.
    \item We provide source-speech-group-held-out and interpreter-disjoint cross-fitting protocols that prevent outer-entity labels from entering upstream quality training, checkpoint selection, or second-stage promptness-model fitting.
    \item We decompose timing, text-derived structural features, and cross-fitted rubric-supervised quality under the same cohort, including a fold-safe residualized-quality audit.
    \item We quantify individual-rater agreement and cross-rater human-label transfer, finding cross-evaluator commonality alongside evaluator-specific calibration; anonymized labels, delay records, rubric, splits, and code are planned for camera-ready release.
\end{itemize}
